\section*{第 \ref{ch:g9:unique-individuals} 章——独特的个体}

\begin{solution}{exo:g9:unique-individuals:1}
减数分裂的洗牌（伴随片段交换）、
\omterm{def:g5:flower-to-fruit:steps}{受精}的抽签，
以及\omterm{prop:g9:genes-and-dna:explains}{突变}带来的逐代渗入。
\end{solution}

\begin{solution}{exo:g9:unique-individuals:2}
它切牌的位置低于整条\omterm{def:g9:chromosomes:chromosomes}{染色体}：同源配对的伙伴
在分离前互换片段，因此即使只有一条\omterm{def:g9:chromosomes:chromosomes}{染色体}
也能传出新的镶嵌组合——$2^{23}$ 这个计数只是下限。
\end{solution}

\begin{solution}{exo:g9:unique-individuals:3}
任意两个人的 \omterm{def:g9:genes-and-dna:dna}{DNA} 文本，
有超过百分之 $99$ 的部分一致；
分散的差异仍使每个人严格独特——亲缘与差异同在一个分子之中。
\end{solution}

\begin{solution}{exo:g9:unique-individuals:4}
它们的生命在共同的起点之后分岔：
获得\omterm{def:g9:heredity-and-traits:traits}{性状}、训练形成的\omterm{def:g8:nervous-communication:synapse}{突触}，以及每个人体内后期发生的\omterm{prop:g9:genes-and-dna:explains}{突变}，
在相同的起点上留下各不相同的记录。
\end{solution}

\begin{solution}{exo:g9:unique-individuals:5}
痕迹来自谁的\omterm{def:g5:first-look-at-cells:cell}{细胞}（用于鉴定），
以及孩子的父母是谁（每个可变位点必须
从两副亲本牌中各发一张）。
\end{solution}

\begin{solution}{exo:g9:unique-individuals:6}
免疫系统能辨识精细的自身标记，
亲缘越近，标记的发牌就越接近——
同一家族的牌共享得最多。
\end{solution}

\begin{solution}{exo:g9:unique-individuals:7}
每位亲本能提供数百万种可能的半副牌；孩子
是母方的一副牌与父方的一副牌的一次配对，
由哪个\omterm{def:g8:sexual-reproduction:gametes}{配子}与哪个\omterm{def:g8:sexual-reproduction:gametes}{配子}相遇决定——两个空间相乘，
其乘积远超人类有史以来出生过的总数，
因此同一副重复的牌实际上从未出现过，
将来也永远不会出现。
\end{solution}

\begin{solution}{exo:g9:unique-individuals:8}
需要说明：克隆个体只是同一副牌的重复——
同一种体质、同一个共同弱点；而由\omterm{def:g2:from-seed-to-plant:seed}{种子}长成的\omterm{def:g8:reproduction-and-environment:population}{种群}，
各不相同的牌会把抗性分散在个体之间，
使部分藤蔓避开害虫的“口味”，
其余植株则落入其口。
\end{solution}

\begin{solution}{exo:g9:unique-individuals:9}
层级：论断从一个\omterm{def:g9:genes-and-dna:gene}{基因}直接跳到家族层面的生活结果，
中间的放大环节被略去。交织：学业成功是一种\omterm{def:g9:heredity-and-traits:heredity}{遗传}范围，
在学校、家庭和
语言中展开——一个\omterm{def:g9:genes-and-dna:gene}{基因}并不等于命运。多样性：它以虚幻的“正常”为标尺衡量分布。
双胞胎：相同的牌
常常导致不同的学业经历，
该标题未通过全部四项检验。
\end{solution}

\begin{solution}{exo:g9:unique-individuals:10}
因为\omterm{def:g5:biodiversity:species}{物种}\emph{就是}其多样性：保留机制把\omterm{def:g9:genes-and-dna:gene}{等位基因}
分散到各个独特个体之中，没有哪一副牌是其他牌
偏离的\omterm{def:g3:sorting-living-things:criterion}{标准}——“正常”所指的是一种没有任何人
真正携带的平均状态，而不是任何人本应
与之相符的文本。
\end{solution}

\begin{solution}{exo:g9:unique-individuals:11}
在每个可变位点上，孩子的两个拼写必须
一个取自母方的一对，一个取自父方的一对。检验
逐个位点进行：这个蓝眼孩子的每一个拼写
都能从这两副牌中发出——略去的部分解释眼睛，
而各个位点
则证明亲子关系。
\end{solution}

\begin{solution}{exo:g9:unique-individuals:12}
举例来说：一次洗牌为每个人从亲本牌中发出一手
从未重复的牌，\omterm{prop:g9:genes-and-dna:explains}{突变}则在每一代中添加前所未见的拼写——
因此人人各不相同。然而每副牌都由同样的四个字母写成，
由同一套密码解读，与任何其他人的牌
每一百个字母中有超过 $99$ 个相同——因此人人同属一家。
这些机制并非彼此对立，而是同一个系统：
共同的文本被不断重新发牌——
亲缘提供牌，
洗牌保证没有任何一手牌重复。
\end{solution}

\begin{solution}{pb:g9:unique-individuals:1}
\textbf{1.} 因为 \omterm{def:g9:genes-and-dna:dna}{DNA} 位于\omterm{def:g6:cell-unit-of-life:parts}{细胞核}中，而
毛干是死亡的角蛋白——毛根的\omterm{def:g5:first-look-at-cells:cell}{细胞}
携带完整解读所需的\omterm{def:g6:cell-unit-of-life:parts}{细胞核}。

\textbf{2.} 可变位点正是个体相异之处——即用于鉴定的分散位点。
共享位点只能证明毛发属于人类，
这是亲缘层面的事实，
对于指认某一个具体的人毫无用处。

\textbf{3.} 这说明毛发\omterm{def:g5:first-look-at-cells:cell}{细胞}与嫌疑人的\omterm{def:g5:first-look-at-cells:cell}{细胞}
共享同一副牌——在有史以来的全部人类中，
只有一个个体会如此，唯一的例外是
同卵双胞胎，他们因同源而共享这副牌。

\textbf{4.} 借助其他证据判定该双胞胎——
是否涉案，依据包括行踪、证人，以及案件中的
“如何到达现场”。证据本身确实无法区分这一对双胞胎，因此
必须由法庭作出判断。

\textbf{5.} 亲属因同源共享大量
可变拼写的比例——多于陌生人，
少于完全匹配：部分重叠程度随
\omterm{prop:g6:classification-kinship:kinship}{亲缘关系}递减，正是发牌过程沿家系树倒推的结果。

\textbf{6.} 水平：图谱是牌组层面的数据——所读取的
可变位点恰恰是为了不反映
\omterm{def:g9:heredity-and-traits:traits}{性状}而挑选的。编织：气质与行为是
在生活中体现出的范围。多样性：“犯罪\omterm{def:g9:genes-and-dna:gene}{基因}”
将人按一个虚幻的正常\omterm{def:g3:sorting-living-things:criterion}{标准}分级。双胞胎：共享同一副牌，
生活各异。图谱指认的是\omterm{def:g5:first-look-at-cells:cell}{细胞}，而非品性。

\textbf{7.} 独立可变位点的数量：
每增加一个位点，都会使符合匹配的人群
进一步缩小——位点足够多时，“全城罕见”
就会变成“全人类罕见”。实验室读取的位点数，
以\omterm{def:g3:sorting-living-things:criterion}{标准}要求为准。

\textbf{8.} 生物学证据断定毛发来自谁的\omterm{def:g5:first-look-at-cells:cell}{细胞}，
但案件仍须确定毛发的来源位置——
何时、以何种方式、是否无辜地出现在现场：匹配
并非叙事，而法庭定罪依据的是
已被证实的叙事。

\textbf{9.} 减数分裂与\omterm{def:g5:flower-to-fruit:steps}{受精}的发牌过程使每副牌
都不可重现，而 \omterm{def:g9:genes-and-dna:dna}{DNA} 的拼写差异使这一
发牌过程可被读取——\omterm{def:g9:chromosomes:chromosomes}{染色体}负责洗牌，
分子则记录下这一过程。

\textbf{10.} 关于 \omterm{def:g9:genes-and-dna:dna}{DNA} 的化学稳定性：文本
既不会在抽屉里褪色，也不会被改写——同一条
双链，历经三十年仍完好，读出的仍是
当初发出的那副牌。

\textbf{11.} 用途：既能排除无辜者，也能确定
亲子关系的身份鉴定；移植配型能够
找到最接近的发牌结果。需警惕的误用：把
\omterm{def:g5:biodiversity:species}{物种}内的多样性读作等级——将差异压缩为
等级，这正是旧时不公背后的错误生物学。

\textbf{12.} 从未有两次\omterm{def:g5:flower-to-fruit:steps}{受精}发出过相同的牌：
双亲各半副牌所构成的可能组合数，乘以
\omterm{def:g8:sexual-reproduction:gametes}{配子}随机组合的结果，超过了有史以来所有人类的总数——
因此全图谱匹配可以指认一个个体，法庭可以
将这根毛发视为“签名”证据。
\end{solution}
